% Generated draft: requires mathematical and source review.
\documentclass[11pt]{article}
\usepackage[margin=25mm]{geometry}
\usepackage{fontspec}
\setmainfont{Latin Modern Roman}
\usepackage{amsmath,amssymb,mathtools,unicode-math}
\setmathfont{Latin Modern Math}
\usepackage{enumitem,needspace,xurl,hyperref}
\setlength{\parindent}{0pt}
\setlength{\parskip}{5pt}
\setlength{\emergencystretch}{3em}
\newcommand{\N}{\mathbb{N}}
\newcommand{\Z}{\mathbb{Z}}
\newcommand{\Q}{\mathbb{Q}}
\newcommand{\R}{\mathbb{R}}
\newcommand{\C}{\mathbb{C}}
\newcommand{\F}{\mathbb{F}}
\newcommand{\A}{\mathbb{A}}
\newcommand{\PP}{\mathbb{P}}
\newcommand{\E}{\mathbb{E}}
\newcommand{\eps}{\varepsilon}
\newcommand{\dd}{\,\mathrm{d}}
\begin{document}
\section*{UnsolvedMath ID: 20000821}
\textbf{Source problem code:} catalogue entry
\subsection{Definitions and mathematical statement}
Let \(C\) be a curve over an algebraically closed field \(k\), and let \(\operatorname{Hilb}^d(C)\), \(d\geq1\), parametrize zero-dimensional closed subschemes of length \(d\). Dimension means the Krull dimension of an irreducible component. The intended nontrivial question is whether a pure one-dimensional curve can have an irreducible Hilbert component \(W\) with
\[
\dim W<d.
\]
The source does not specify purity, reducedness, irreducibility, or characteristic. Without a pure-curve convention, isolated zero-dimensional components would change the question. For an integral curve the open locus of \(d\) distinct smooth points has dimension \(d\); its closure is the smoothable component. A subscheme is smoothable here if it is a flat limit, on the same curve, of such reduced smooth-supported configurations.

The source remark reports a nonsmoothable component of dimension exactly \(d>1\). This does not answer whether a component of smaller dimension exists. The remark gives neither its curve nor a reference, and nonsmoothable must not be equated with smaller-dimensional.

The attached aid supplies a support-separated reduction for integral curves with isolated singularities \(p_1,\ldots,p_s\). Fix local lengths \(a_j\geq0\), \(\sum_j a_j\leq d\), and chosen punctual families \(Q_j\) of subschemes supported at \(p_j\), with dimensions \(q_j\). Let the remaining length \(d-\sum_j a_j\) vary over the smooth locus. The reported stratum dimension is
\[
d-\sum_j a_j+\sum_j q_j.
\]
For its closure to be a counterexample component, both conditions are required: a genuine punctual deficit \(\sum_jq_j<\sum_ja_j\), and global maximality of that closure among irreducible closed subsets of \(\operatorname{Hilb}^d(C)\). A small punctual family can lie inside the closure of a larger deforming family and hence need not be a component.

For irreducible curves with one rational \(r\)-fold singularity, the aid invokes a 2025 classification and reports that the nonsmoothable components indexed by an admissible \(a\) have dimension
\[
D_{r,a}(d)=d+(a-1)(r-a)-1.
\]
It says every admitted component has dimension at least \(d\), with equality only at \((r,a)=(3,2)\). The local singularity convention and complete admissible \(a,d\) range are not given; this formula must not be evaluated at arbitrary indices and treated as a component-existence assertion. In particular this test family supplies no smaller-dimensional component. The aid describes the unrestricted pure-curve question as still open in the literature it checked. Its reduction and classification application are supplied unreviewed partial claims.
\subsection{Short English statement}
Can the length-\(d\) Hilbert scheme of a pure curve have an irreducible component smaller than \(d\)? The supplied aid separates small support strata from actual components and reports that one rational-fold-singularity family has no such deficit, despite possible nonsmoothable components of dimension \(d\).
\subsection{Sources}
\begin{itemize}
\item AIM Problem 20, \emph{Components of Hilbert schemes}, including the dimension-\(d\) nonsmoothable-component remark: \url{https://aimath.org/WWN/hilbertschemes/hilbertschemes.pdf}.
\item Local registry \texttt{data/unsolvedmath/problems.json}, permanent ID \(20000821\), exact source code \texttt{catalogue entry}; canonical \texttt{aim-arithmetic-geometry-notes.json}, notes[66].
\item The attached aid invokes a 2025 classification of curves with a rational \(r\)-fold singularity, without authors, title, identifier, or full statement. Attempt 1 is explicitly machine-generated, not peer reviewed, and requires independent verification before citation.
\end{itemize}
\end{document}
